% Generated draft: requires mathematical and source review.
\documentclass[11pt]{article}
\usepackage[margin=25mm]{geometry}
\usepackage{fontspec}
\setmainfont{Latin Modern Roman}
\usepackage{amsmath,amssymb,mathtools,unicode-math}
\setmathfont{Latin Modern Math}
\usepackage{enumitem,needspace,xurl,hyperref}
\setlength{\parindent}{0pt}
\setlength{\parskip}{5pt}
\setlength{\emergencystretch}{3em}
\newcommand{\N}{\mathbb{N}}
\newcommand{\Z}{\mathbb{Z}}
\newcommand{\Q}{\mathbb{Q}}
\newcommand{\R}{\mathbb{R}}
\newcommand{\C}{\mathbb{C}}
\newcommand{\F}{\mathbb{F}}
\newcommand{\A}{\mathbb{A}}
\newcommand{\PP}{\mathbb{P}}
\newcommand{\E}{\mathbb{E}}
\newcommand{\eps}{\varepsilon}
\newcommand{\dd}{\,\mathrm{d}}
\begin{document}
\section*{UnsolvedMath ID: 30006342}
\textbf{Source problem code:} catalogue entry
\subsection{Definitions and mathematical statement}
Let \(\Omega\) be a finite set and let \(G\leq\operatorname{Sym}(\Omega)\) be a primitive permutation group. Here primitive means transitive with no \(G\)-invariant partition other than the partition into singletons and the one-part partition. For \(B\subseteq\Omega\), its pointwise stabilizer is
\[
G_{(B)}=\{g\in G:g(\beta)=\beta\text{ for every }\beta\in B\}.
\]
A base is a subset \(B\subseteq\Omega\) with \(G_{(B)}=\{1\}\), and the base size is
\[
b(G)=\min\{|B|:B\subseteq\Omega,\ G_{(B)}=\{1\}\}.
\]
Assume \(b(G)\geq2\). Define the generalized Saxl graph \(\Sigma(G)\) to be the simple undirected graph with vertex set \(\Omega\) in which distinct vertices \(\alpha,\beta\) are adjacent precisely when there exists a base \(B\) of cardinality \(b(G)\) containing both vertices. Write \(N_{\Sigma(G)}(\alpha)\) for the set of neighbors of \(\alpha\).

\textbf{Common-Neighbor Conjecture for Generalized Saxl Graphs.} Is it true that, for every such \(G\),
\[
\forall\alpha,\beta\in\Omega,\quad
\alpha\neq\beta\ \Longrightarrow\
N_{\Sigma(G)}(\alpha)\cap N_{\Sigma(G)}(\beta)\neq\varnothing?
\]
Equivalently, for every distinct \(\alpha,\beta\in\Omega\), must there exist \(\gamma\in\Omega\setminus\{\alpha,\beta\}\) and bases \(B_1,B_2\subseteq\Omega\) satisfying
\[
|B_1|=|B_2|=b(G),\qquad
\{\alpha,\gamma\}\subseteq B_1,\qquad
\{\beta,\gamma\}\subseteq B_2?
\]
The two bases need not coincide. When \(b(G)=2\), this definition gives the usual Saxl graph. The phrase ``every two vertices'' is interpreted above as distinct vertices. Allowing equal vertices gives an equivalent formulation here: transitivity ensures that every vertex belongs to a minimum base, and \(b(G)\geq2\) therefore ensures that every vertex has a neighbor.

The question is reconstructed from the supplied statement and clean statement. The supplied original excerpt refers to a conjecture without stating it, so that excerpt alone does not independently establish this formulation. The supplied dated curation labels the problem partially solved and reports results for several families, including many sporadic almost simple groups; these are source qualifications, not independently verified conclusions of this draft.
\subsection{Short English statement}
For every finite primitive permutation group with minimum base size at least two, join two distinct points when they occur together in a minimum base. The question asks whether every pair of distinct points has a common neighbor. The minimum bases witnessing the two adjacencies may differ. The supplied record reports partial results, but the general question and its literature status require human review.
\subsection{Sources}
\begin{itemize}
\item Supplied local UnsolvedMath registry record: permanent ID \(30006342\), title ``Common-Neighbor Conjecture for Generalized Saxl Graphs,'' exact source problem code \texttt{catalogue entry}, in the Oberwolfach Reports Open Problems collection.
\item Supplied report citation: \emph{Computational Group Theory} (2025), Oberwolfach Reports, pp.~1389--1450, DOI \texttt{10.4171/owr/2025/27}. Source URL and literature-source label ``OWR report'': \url{https://doi.org/10.4171/owr/2025/27}.
\item Attribution supplied in the original excerpt: ``Bases for primitive sporadic almost simple permutation groups,'' Saul D. Freedman, joint work with Hong Yi Huang, Melissa Lee, and Kamilla Rekv\'enyi. The record does not identify a proposer of the conjecture.
\item Supplied literature source, ``On generalized Saxl graphs'': \url{https://arxiv.org/abs/2410.22613}.
\item Supplied literature source, ``Published generalized Saxl graphs'': \url{https://doi.org/10.5802/alco.493}.
\item Literature assessment dated exactly \texttt{2026-08-22}: supplied dated curation, not independently verified. It states that the Common Neighbour Conjecture for Saxl and generalized Saxl graphs remains open in general and that the primary work proves cases in several families, including many sporadic almost simple groups. The background expressly calls this dated literature triage derived from the supplied Oberwolfach curation and requests independent verification of time-sensitive status and cited claims.
\item The supplied registry labels the statement \texttt{corrected\_verified} and reports that the official Oberwolfach report was checked, the question isolated from surrounding exposition, and notation repaired. This records the supplied verification claim; no source visit or independent verification is asserted here.
\end{itemize}
\end{document}
